\documentclass{article}
\usepackage[italian]{babel}
\usepackage[utf8]{inputenc}
\usepackage[T1]{fontenc}
\usepackage{graphicx}
\usepackage[colorinlistoftodos]{todonotes}
\usepackage[colorlinks=true, allcolors=tudelftblue]{hyperref} %sets hyperlink colour
\usepackage{caption}
\usepackage{subcaption}
\usepackage{xcolor}
\usepackage{roboto} % for Roboto Slab font
\usepackage{float}
\usepackage{titling}
\usepackage{blindtext}\usepackage{titlesec}
\usepackage[square,sort,comma,numbers]{natbib}
\usepackage[colorinlistoftodos]{todonotes}
\usepackage{tikz}
\usepackage{geometry}
\usepackage{sectsty}
\usepackage{amsmath}
\usepackage{tikzpagenodes}
\usepackage{booktabs}
\usepackage{listings}
%%% Pacchetti aggiunti per il progetto
\usepackage{tabularx}   % glossari a larghezza fissa
\usetikzlibrary{shapes.geometric} % rombi degli schemi scheletro
\definecolor{tudelftdarkblue}{RGB}{0,0,0}
\definecolor{tudelftcyan}{RGB}{209,65,36}
\definecolor{tudelftblue}{RGB}{99, 102, 106}
\geometry{a4paper, margin=2cm}
\setlength{\parindent}{0pt}
\allsectionsfont{\color{black}} %sets colour for all headers
\usepackage{helvet}
\renewcommand{\familydefault}{\sfdefault}
\sectionfont{\fontfamily{RobotoSlab-TLF}\selectfont}
\captionsetup{font=small}
% Paragrafi dei requisiti piu' compatti (ridefinito dopo sectsty, che altrimenti impone 3.25ex)
\makeatletter
\renewcommand\paragraph{\@startsection{paragraph}{4}{\z@}{1.2ex \@plus .4ex \@minus .2ex}{-0.6em}{\normalfont\normalsize\bfseries\color{black}}}
\makeatother
%%%%%%%%%%%%%%%%%%%%%%%%%%%%%%%%%%%%%%%%%%%%%%%%%%%%%%%%%
%%% Macro del progetto
% Concetto del glossario della vista (corsivo nel testo dei requisiti)
\newcommand{\concetto}[1]{\textit{#1}}
% Colori dei costrutti evidenziati: IO = instance-of, PO = part-of, G = gerarchia
\definecolor{colIO}{RGB}{0,90,170}
\definecolor{colPO}{RGB}{0,125,60}
\definecolor{colG}{RGB}{209,65,36}
% Etichetta del costrutto (#1 = IO, PO oppure G)
\newcommand{\etichetta}[1]{{\setlength{\fboxsep}{0.8pt}\colorbox{col#1}{\textcolor{white}{\scriptsize\bfseries #1}}}}
% Frase che esprime il costrutto, nel colore del costrutto e seguita dall'etichetta
\newcommand{\evid}[2]{\textcolor{col#1}{\textbf{#2}}\,\etichetta{#1}}
% Colonna di testo allineata a sinistra per tabularx
\newcolumntype{Y}{>{\raggedright\arraybackslash}X}
% Stili degli schemi scheletro: rettangolo = entita', rombo = relazione
\tikzset{
  entita/.style={draw, semithick, rectangle, minimum width=2.5cm, minimum height=0.8cm, font=\small},
  relazione/.style={draw, semithick, diamond, aspect=2.4, inner sep=1.5pt, font=\footnotesize},
  every path/.style={semithick},
}
%%%%%%%%%%%%%%%%%%%%%%%%%%%%%%%%%%%%%%%%%%%%%%%%%%%%%%%%%
\begin{document}

\begin{titlepage}
    \fontfamily{RobotoSlab-TLF}\selectfont 
%%%%%%%%%%%%%%%%%%%%%%%%%%%%%%%%%%%%%%%%%%%%%%%%%%%%%%%%%%UNCOMMENT THE FOLLOWING FOR LESS "PLAIN" TITLE PAGE (SELECT WITH MOUSE AND PRESS CTRL AND /)

    % \begin{tikzpicture}[remember picture,overlay]
    %     % Set seed for random number generator
    %     \pgfmathsetseed{4}
    %     % Define the text area to avoid
    %     \path (current page text area.south west) rectangle (current page text area.north east);
    %     % Adding circles spread over the entire page
    %     \foreach \x in {1,...,1000}
    %         \draw[tudelftdarkblue] (current page.south west) ++(rand*\paperwidth,rand*\paperheight) circle (rand*0.3);
    %     % Define coordinates for the corners of the white box
    %     \coordinate (A) at ([shift={(-8cm,12cm)}]current page.center);
    %     \coordinate (B) at ([shift={(8cm,-5cm)}]current page.center);
    %     % Draw the white background box
    %     \fill[white] (A) rectangle (B);
    %     % Adding equations as background features
    %     \node[anchor=center,rotate=20,text=tudelftcyan] at ([shift={(-7cm,-2cm)}]current page.center) {\fontsize{18}{22}\selectfont
    %     $\nabla^2 T - \frac{1}{\alpha}\frac{\partial T}{\partial t} = 0$};
    %     \node[anchor=center,rotate=-15,text=tudelftcyan] at ([shift={(5cm,-4cm)}]current page.center) {\fontsize{18}{22}\selectfont
    %     $\frac{\partial \rho}{\partial t} + \nabla \cdot (\rho \mathbf{v}) = 0$};
    %     \node[anchor=center,rotate=20,text=tudelftcyan] at ([shift={(6cm,4cm)}]current page.center) {\fontsize{18}{22}\selectfont
    %     $a^2 + b^2 = c^2$};
    %     \node[anchor=center,rotate=10,text=tudelftcyan] at ([shift={(7cm,-2cm)}]current page.center) {\fontsize{18}{22}\selectfont
    %     $E = \frac{\sigma}{\varepsilon}$};
    %     \node[anchor=center,rotate=-10,text=tudelftcyan] at ([shift={(-6cm,4cm)}]current page.center) {\fontsize{18}{22}\selectfont
    %     $F = ma$};
    %     \node[anchor=center,rotate=5,text=tudelftcyan] at ([shift={(-4cm,-5cm)}]current page.center) {\fontsize{18}{22}\selectfont
    %     $Q = -\frac{kA}{\mu} \frac{\Delta P}{L}$};
    % \end{tikzpicture}
%%%%%%%%%%%%%%%%%%%%%%%%%%%%%%%%%%%%%%%%%%%%%%%%%%%%%%%%%%
    \vspace*{3cm}
    
    \centering
    {\Huge \textbf{\textcolor{black}{System Logistics - Copper Golem}}}\\[0.4cm]
    {\Large \textbf{Gestione di un impianto intralogistico}}\\[1.5cm]
    \textsc{\LARGE Basi di Dati}\\[0.5cm]
    \text{\large MN1-525}\\[2cm]

    {\Large \textbf{\textcolor{tudelftdarkblue}{Autori:}}}\\[0.5cm]
    \begin{tabular}{c}
        \Large \textcolor{tudelftdarkblue}{Gottardi Emanuele (210373)} \\
        \Large \textcolor{tudelftdarkblue}{Rossato Emanuele (210372)} \\
        \Large \textcolor{tudelftdarkblue}{Soncini Leonardo (210382)} \\
    \end{tabular}\\[2cm]
    
    {\Large \textcolor{tudelftdarkblue}{\today}}
    
    \vfill
    \begin{center}
        \includegraphics[width=0.6\textwidth]{images/Logo_C_Positivo_Colore.png}
    \end{center}
\end{titlepage}

%%% Create a table of contents
\tableofcontents
\newpage

\addcontentsline{toc}{section}{Introduzione}
\section*{Introduzione}
System Logistics S.p.A.~\cite{sl_web} presenta come core business la creazione di impianti intralogistici automatici, per settori alimentari, come fruizione di cibo e bevande e per la grande distribuzione. 
Un tipico sistema prodotto da loro comprende magazzini automatici serviti da traslo-elevatori, sistemi di trasporto (come trasportatori a rulli e a catene), veicoli a guida automatica e stazioni di picking, governati da un software centrale chiamato SYSTORE. 
Una volta avviato il sistema, l’azienda fornisce assistenza sotto forma di help desk, manutenzione e fornitura di ricambi.

Il progetto riguarderà la base di dati che SYSTORE comprende relativa ad un singolo impianto in uno stabilimento del committente, in questo caso un produttore di bevande. 
Il database fornirà supporto all’intera area logistica, comprendendo una totale copertura dello stoccaggio del magazzino e dell’invio della merce.
Saranno quindi considerate tre classi di utenti distinti: il responsabile delle operazioni di magazzino (vista A), il tecnico dell’assistenza fornita da System Logistics (vista B) e infine l’analista delle prestazioni e del gemello digitale (vista C).






%=========================================================
\section{Definizione dei requisiti}
\label{sec:requisiti}
\textbf{Legenda.} Nel testo dei requisiti il \concetto{corsivo} indica i concetti del glossario della
vista; le frasi evidenziate indicano \evid{IO}{relazioni instance-of}, \evid{PO}{relazioni part-of} e
\evid{G}{gerarchie}, come definite in~\cite{atzeni2018}.

%---------------------------------------------------------
\subsection{Vista A -- Responsabile delle operazioni di magazzino}
\label{sec:req-A}
Il responsabile delle operazioni di magazzino gestisce i flussi fisici della merce, ed è un dipendente del committente che lavora con il Warehouse Management System del software SYSTORE. Registra l’arrivo delle unità di carico dalla produzione, consulta le giacenze di ogni articolo e ubicazione, blocca e sblocca i lotti (individuando quelli con la scadenza più vicina), rilascia gli ordini di spedizione e controlla le missioni in corso su ogni macchina.
Per compiere ciò, per ogni articolo si registrano codice, descrizione, formato di vendita, elementi per pallet e durata di conservazione. La merce arriva in lotti di produzione: ogni lotto è la controparte fisica di un articolo ed è identificato dal codice dell’articolo stesso e da un secondo identificatore numerico, univoco all’interno dello stesso prodotto. Presenta inoltre la data di produzione e di scadenza. 
La merce è spostata e sistemata su unità di carico, le quali si dividono in pallet e contenitori, questi ultimi presenti solo nelle aree di picking. Ogni unità ha un codice univoco, un tipo, un’altezza e un peso misurati al ricevimento. Solitamente le unità contengono un solo tipo di articolo, ma in fase di picking è possibile comporre un pallet inserendone più di uno; per ogni articolo selezionato si registrano lotto e numero di elementi.
Un’ubicazione è una posizione che può ospitare una unità di carico, con un codice univoco e uno stato (libera, occupata o bloccata). Le ubicazioni si dividono in celle e baie. Il magazzino è diviso in corridoi, ognuno servito da un traslo-elevatore, e ogni corridoio è composto da celle: una cella è individuata nel corridoio da lato, colonna, livello e profondità, ha altezza e portata massime e ospita al più una unità di carico. Le baie sono i punti di scambio con l’esterno del magazzino e hanno un tipo, che ne specifica il compito rispetto al magazzino (come ingresso, uscita, picking o carico). 
Ogni spostamento di una unità da un’ubicazione di origine a una di destinazione è una missione, creata dal WMS con codice, priorità, stato (in coda, in esecuzione, completata o annullata) e istanti di creazione, inizio e fine. Le missioni si distinguono in missioni di ingresso, di uscita e di trasferimento: le prime portano una UDC da una baia a una cella e registrano l’esito del controllo di sagoma e peso, per assicurarsi la correttezza del prodotto; le seconde la portano da una cella a una baia per evadere un ordine; le terze la spostano tra due celle o tra due baie e ne registrano il motivo. Una missione impiega di solito più macchine in sequenza e per ognuna si registrano la posizione nella sequenza e gli istanti di inizio e fine del tratto.
Le missioni sono eseguite dalle macchine dell’impianto: ognuna di esse può essere la rappresentazione di un modello scelto tra traslo-elevatori, trasportatori e veicoli laser comandati. Di ogni macchina si consultano codice, descrizione e stato (in lavoro, in attesa, in errore, in manutenzione). 
Gli ordini di spedizione arrivano dal gestionale del committente verso SYSTORE, ciascuno con numero, cliente destinatario, data di spedizione prevista, baia di carico e stato, e con il numero di elementi richiesto per ogni articolo. 

%---------------------------------------------------------
\subsection{Vista B -- Tecnico dell'assistenza (manutenzione e ricambi)}
\label{sec:req-B}

Un tecnico di System Logistics riceve tramite contatto del cliente richieste di assistenza o di
manutenzione presso la loro sede, una volta accolta la richiesta si collega in remoto presso il
Systore installato dal cliente e controlla prima stato, avvisi e modelli installati. Ogni volta che
agisce in loco presso le macchine del cliente apre un intervento, dove inserisce ore lavorate,
nominativi di chi ha svolto il lavoro ed eventualmente ricambi utilizzati. Ogni volta che torna dal
cliente chiude l'intervento indipendentemente se sia finito o se debba tornarci il giorno dopo.
Inoltre Systore avvisa System Logistics di interventi preventivi in scadenza da programmare, oppure
se alcuni componenti di facile usura sono vicini al fine vita utile.
Ogni tecnico può inoltre consultare lo storico dei guasti e degli interventi effettuati su ogni
macchina o controllare i ricambi sotto scorta.

\paragraph{Macchina.}
Di ogni macchina installata si registrano matricola, anno di installazione, posizione nell'impianto
all'interno dello stabile del cliente e contatori di utilizzo (ore di funzionamento e cicli). Ogni
macchina è un esemplare di un tipo di macchina (modello) del catalogo System Logistics, con codice,
categoria (traslo-elevatore, trasportatore, LGV), portata e velocità nominali.

\paragraph{Componente.}
Per ogni tipo è definito il piano di manutenzione preventiva: un elenco di attività, ognuna con
descrizione e periodicità (in ore di funzionamento o in giorni). Ogni macchina è composta da
componenti sostituibili (es.\ motore, gruppo forche, cinghie o cuscinetti). Ogni componente è
individuato nella macchina da un codice di posizione e ha data di montaggio, ore di funzionamento e
solo se prevista una vita utile; quando le ore si avvicinano alla vita utile il sistema ne propone la
sostituzione.

\paragraph{Guasto.}
Quando una macchina si ferma per un errore, il PLC invia un allarme che il sistema registra come
guasto, con istante, codice di errore, descrizione e gravità. Un guasto resta aperto finché non è
risolto, ma non è detto serva un tecnico per risolverlo.
Ogni intervento riguarda una macchina o un gruppo di macchine ed ha numero, data, orari di inizio e
fine, descrizione del lavoro ed esito. Gli interventi si distinguono in preventivi e correttivi: i
primi svolgono le attività del piano giunte a scadenza; i secondi risolvono uno o più guasti e
registrano la causa individuata e la durata del fermo.
Di ogni ricambio conservato nel deposito si registrano codice, descrizione, quantità disponibile,
scorta minima e componenti che può sostituire (in quanto un componente può sostituire più componenti
con poche modifiche). Per ogni intervento si registrano i ricambi prelevati e le quantità; sotto la
scorta minima il tecnico chiede il riordino alla sede.

\paragraph{Tecnico.}
Di ogni tecnico si registrano matricola, nome, cognome, telefono e specializzazione (meccanica,
elettrica, software). A un intervento partecipano uno o più tecnici, ciascuno con le ore lavorate,
in quanto possono arrivare in momenti diversi della giornata.

%---------------------------------------------------------
\subsection{Vista C -- Analista delle prestazioni e del gemello digitale}
\label{sec:req-C}
L'analista lavora come gestore di progetto relativo ad un singolo committente di System Logistics e segue da remoto il rendimento
dell'impianto dopo l'avviamento~\cite{sl_cs}. Deve capire se le \concetto{macchine} lavorano come
dovrebbero e, per farlo, mette a confronto i dati raccolti sul campo con le previsioni del
\concetto{gemello digitale} (la copia virtuale dell'impianto su cui si fanno le
simulazioni). Consulta le misure dei \concetto{sensori} e gli \concetto{eventi} di
ogni \concetto{macchina}, controlla gli \concetto{indicatori} calcolati a fine turno e ricalibra il
\concetto{gemello digitale} quando le sue previsioni si allontanano dalla realtà. Prepara inoltre gli
\concetto{scenari} da simulare, confronta valori reali e previsti e segnala al tecnico dell'assistenza
le \concetto{macchine} che rendono meno del previsto. Le sue simulazioni servono anche al committente,
che prima di un periodo di picco vuole sapere se l'impianto reggerà.

\paragraph{Sensori e misure.}
Gran parte dei dati di campo arriva dai \concetto{sensori}. Di ognuno si registrano un codice, la
grandezza misurata (per esempio la corrente assorbita da un motore o una temperatura), l'unità di misura
e ogni quanto viene letto. \evid{G}{Un \concetto{sensore} può essere di macchina, se è montato su una
\concetto{macchina}, oppure ambientale}: in questo caso controlla una zona dell'impianto e ha un
intervallo di valori ammessi, come i termometri dell'area di stoccaggio, dove le bevande non devono
scaldarsi troppo. Ogni lettura di un \concetto{sensore} produce una misura, con istante e valore.

\paragraph{Eventi.}
Oltre alle misure, le \concetto{macchine} comunicano degli \concetto{eventi}: li invia il controllore di
bordo (PLC) attraverso SYSTORE, ognuno con il proprio istante. \evid{G}{Gli \concetto{eventi} si dividono in cambi di stato, allarmi e fine missione}. Un
cambio di stato registra lo stato lasciato e quello nuovo (in lavoro, in attesa, in errore, in
manutenzione). Un allarme ha un codice e una gravità, che va dal semplice avviso fino al blocco della
\concetto{macchina}. Un \concetto{evento} di fine missione indica la missione che la
\concetto{macchina} ha appena concluso.

\paragraph{Indicatori.}
Gli \concetto{indicatori} di prestazione (KPI) sono definiti una volta sola. Ciascuno con
codice, nome, formula di calcolo e unità di misura. Si usano la disponibilità, il throughput (UDC
movimentate all'ora), il tempo medio di ciclo e l'energia consumata per missione. Alla fine di ogni
turno, individuato da data e fascia oraria, il sistema calcola il valore reale di ogni
\concetto{indicatore} per ogni \concetto{macchina}: la disponibilità deriva dalla durata degli stati,
throughput e tempo di ciclo dalle missioni concluse, l'energia dalle misure dei \concetto{sensori}. Il
termine di paragone è il valore previsto dal \concetto{gemello digitale}: ogni \concetto{indicatore} ha
una soglia che dice di quanto il valore reale può allontanarsene prima che la \concetto{macchina} venga
segnalata. Se per esempio la soglia della disponibilità è di tre punti e la previsione è del 96\%, un
valore reale del 92\% fa scattare la segnalazione. Alcuni \concetto{indicatori}, come la disponibilità,
hanno anche un valore obiettivo fissato nel contratto di assistenza.

\paragraph{Gemello digitale.}
Il \concetto{gemello digitale} riproduce in simulazione le \concetto{macchine} dell'impianto e il loro
funzionamento. Ogni volta che viene ricalibrato sui dati di campo ne nasce una nuova versione, con numero
progressivo, data di calibrazione ed errore medio delle previsioni rispetto ai valori reali.
\evid{PO}{Ogni versione del \concetto{gemello digitale} è composta dalle repliche virtuali delle singole
\concetto{macchine}}: una replica corrisponde a una \concetto{macchina} e ne contiene i parametri di
funzionamento (velocità, accelerazioni, tempi di presa e deposito), che possono cambiare da una versione
all'altra. Per ogni versione si indica anche il livello di integrazione con
l'impianto~\cite{kritzinger2018}: modello digitale se i dati sono scambiati a mano, ombra digitale se
arrivano in automatico dall'impianto alla simulazione, gemello vero e proprio se viaggiano in automatico
nei due sensi.

\paragraph{Scenari e simulazioni.}
Uno \concetto{scenario} stabilisce le condizioni in cui far lavorare il \concetto{gemello digitale}. Ha
un nome, una descrizione e le condizioni ipotizzate, come un diverso profilo degli ordini o la modifica
di alcune repliche, per simulare il fermo di una \concetto{macchina} o una velocità maggiore. Lo
\concetto{scenario} di riferimento riproduce le condizioni effettive di ogni turno e serve al confronto
quotidiano: dagli scarti tra le sue previsioni e i valori reali si ricava anche l'errore medio di ogni
versione. Gli altri \concetto{scenari} sono ipotesi del tipo «cosa succede se», come un picco di ordini
prima dell'estate o l'aggiunta di un turno notturno. \evid{IO}{Ogni simulazione è un'esecuzione
concreta di uno \concetto{scenario}} su una versione del \concetto{gemello digitale}, con data e durata,
e restituisce il valore previsto di ogni \concetto{indicatore} per ogni \concetto{macchina}. La
simulazione resta legata alla versione usata, così si possono confrontare previsioni ottenute con
calibrazioni diverse.


%=========================================================
\section{Analisi dei requisiti e schema scheletro}
\label{sec:analisi}

%---------------------------------------------------------
\subsection{Vista A -- Glossario e schema scheletro}
\label{sec:an-A}

\begin{table}[H]
\centering\small
\begin{tabularx}{\textwidth}{@{}>{\raggedright\arraybackslash}p{2.0cm} Y >{\raggedright\arraybackslash}p{3.0cm} >{\raggedright\arraybackslash}p{5.2cm}@{}}
\toprule
\textbf{Termine} & \textbf{Descrizione} & \textbf{Sinonimi} & \textbf{Collegamenti} \\
\midrule
Articolo   & Prodotto del committente        & --                      & UDC, Ordine \\
UDC        & Unità fisica che porta la merce & Unità di carico, pallet & Articolo, Ubicazione, Missione \\
Ubicazione & Posto che può ospitare una UDC  & Posizione               & UDC, Missione \\
Missione   & Trasporto di una UDC            & Spostamento             & UDC, Ubicazione, Macchina, Ordine \\
Macchina   & Dispositivo che esegue missioni & --                      & Missione \\
Ordine     & Richiesta di spedizione         & Ordine di spedizione    & Articolo, Missione \\
\bottomrule
\end{tabularx}
\caption{Glossario dei concetti della vista A.}
\label{tab:gloss-A}
\end{table}

\begin{figure}[H]
\centering
\begin{tikzpicture}
  % entita'
  \node[entita] (art) at (0,2.8)    {Articolo};
  \node[entita] (ord) at (6.8,2.8)  {Ordine};
  \node[entita] (udc) at (0,0)      {UDC};
  \node[entita] (mis) at (6.8,0)    {Missione};
  \node[entita] (mac) at (13.6,0)   {Macchina};
  \node[entita] (ubi) at (3.4,-2.8) {Ubicazione};
  % relazioni
  \node[relazione] (ric) at (3.4,2.8)   {Richiesta};
  \node[relazione] (con) at (0,1.4)     {Contenuto};
  \node[relazione] (eva) at (6.8,1.4)   {Evasione};
  \node[relazione] (mov) at (3.4,0)     {Movimentazione};
  \node[relazione] (ese) at (10.2,0)    {Esecuzione};
  \node[relazione] (col) at (0,-1.4)    {Collocazione};
  \node[relazione] (ori) at (6.0,-1.05) {Origine};
  \node[relazione] (des) at (7.8,-2.0)  {Destinazione};
  % collegamenti
  \draw (art) -- (ric) -- (ord);
  \draw (art) -- (con) -- (udc);
  \draw (ord) -- (eva) -- (mis);
  \draw (udc) -- (mov) -- (mis);
  \draw (mis) -- (ese) -- (mac);
  \draw (udc) -- (col) |- (ubi.west);
  \draw ([xshift=-0.8cm]mis.south) -- (ori) |- ([yshift=0.2cm]ubi.east);
  \draw ([xshift=1.0cm]mis.south) -- (des) |- ([yshift=-0.2cm]ubi.east);
\end{tikzpicture}
\caption{Schema scheletro della vista A.}
\label{fig:sch-A}
\end{figure}

%---------------------------------------------------------
\subsection{Vista B -- Glossario e schema scheletro}
\label{sec:an-B}

\begin{table}[H]
\centering\small
\begin{tabularx}{\textwidth}{@{}>{\raggedright\arraybackslash}p{2.0cm} Y >{\raggedright\arraybackslash}p{3.4cm} >{\raggedright\arraybackslash}p{5.4cm}@{}}
\toprule
\textbf{Termine} & \textbf{Descrizione} & \textbf{Sinonimi} & \textbf{Collegamenti} \\
\midrule
Macchina   & Esemplare in servizio nel sito  & Apparecchiatura   & Componente, Guasto, Intervento \\
Componente & Parte sostituibile di macchina  & Gruppo, organo    & Macchina, Ricambio \\
Guasto     & Malfunzionamento di macchina    & Allarme, errore   & Macchina, Intervento \\
Intervento & Lavoro svolto su una macchina   & Ordine di lavoro  & Macchina, Guasto, Tecnico, Ricambio \\
Tecnico    & Addetto all'assistenza tecnica  & Manutentore       & Intervento \\
Ricambio   & Pezzo a scorta nel deposito     & Pezzo di ricambio & Componente, Intervento \\
\bottomrule
\end{tabularx}
\caption{Glossario dei concetti della vista B.}
\label{tab:gloss-B}
\end{table}

\begin{figure}[H]
\centering
\begin{tikzpicture}
  % entita'
  \node[entita] (gua) at (0,2.8)    {Guasto};
  \node[entita] (mac) at (6.8,2.8)  {Macchina};
  \node[entita] (com) at (13.6,2.8) {Componente};
  \node[entita] (int) at (6.8,0)    {Intervento};
  \node[entita] (ric) at (13.6,0)   {Ricambio};
  \node[entita] (tec) at (6.8,-2.8) {Tecnico};
  % relazioni
  \node[relazione] (seg) at (3.4,2.8)  {Segnalazione};
  \node[relazione] (cmp) at (10.2,2.8) {Composizione};
  \node[relazione] (man) at (6.8,1.4)  {Manutenzione};
  \node[relazione] (cpt) at (13.6,1.4) {Compatibilità};
  \node[relazione] (ris) at (3.4,0)    {Risoluzione};
  \node[relazione] (uti) at (10.2,0)   {Utilizzo};
  \node[relazione] (par) at (6.8,-1.4) {Partecipazione};
  % collegamenti
  \draw (gua) -- (seg) -- (mac);
  \draw (mac) -- (cmp) -- (com);
  \draw (mac) -- (man) -- (int);
  \draw (com) -- (cpt) -- (ric);
  \draw (gua) |- (ris) -- (int);
  \draw (int) -- (uti) -- (ric);
  \draw (int) -- (par) -- (tec);
\end{tikzpicture}
\caption{Schema scheletro della vista B.}
\label{fig:sch-B}
\end{figure}

%---------------------------------------------------------
\subsection{Vista C -- Glossario e schema scheletro}
\label{sec:an-C}

\begin{table}[H]
\centering\small
\begin{tabularx}{\textwidth}{@{}>{\raggedright\arraybackslash}p{2.0cm} Y >{\raggedright\arraybackslash}p{3.4cm} >{\raggedright\arraybackslash}p{5.4cm}@{}}
\toprule
\textbf{Termine} & \textbf{Descrizione} & \textbf{Sinonimi} & \textbf{Collegamenti} \\
\midrule
Macchina         & Fonte di eventi e misure          & --                     & Sensore, Evento, Indicatore, Gemello digitale \\
Sensore          & Dispositivo di misura             & Trasduttore            & Macchina \\
Evento           & Fatto comunicato da una macchina  & Notifica, segnale      & Macchina \\
Indicatore       & Misura sintetica di prestazione   & KPI, prestazione       & Macchina, Scenario \\
Scenario         & Condizioni ipotetiche da simulare & Ipotesi, caso what-if  & Indicatore, Gemello digitale \\
Gemello digitale & Replica virtuale dell'impianto    & Modello di simulazione & Macchina, Scenario \\
\bottomrule
\end{tabularx}
\caption{Glossario dei concetti della vista C.}
\label{tab:gloss-C}
\end{table}

\begin{figure}[H]
\centering
\begin{tikzpicture}
  % entita'
  \node[entita] (eve) at (6.8,2.8)   {Evento};
  \node[entita] (sen) at (0,0)       {Sensore};
  \node[entita] (mac) at (6.8,0)     {Macchina};
  \node[entita, minimum width=2.9cm] (gem) at (13.6,0) {Gemello digitale};
  \node[entita] (ind) at (6.8,-2.8)  {Indicatore};
  \node[entita] (sce) at (13.6,-2.8) {Scenario};
  % relazioni
  \node[relazione] (gen) at (6.8,1.4)   {Generazione};
  \node[relazione] (ins) at (3.4,0)     {Installazione};
  \node[relazione] (rap) at (10.2,0)    {Rappresentazione};
  \node[relazione] (cal) at (6.8,-1.4)  {Calcolo};
  \node[relazione] (sim) at (13.6,-1.4) {Simulazione};
  \node[relazione] (pre) at (10.2,-2.8) {Previsione};
  % collegamenti
  \draw (eve) -- (gen) -- (mac);
  \draw (sen) -- (ins) -- (mac);
  \draw (mac) -- (rap) -- (gem);
  \draw (mac) -- (cal) -- (ind);
  \draw (gem) -- (sim) -- (sce);
  \draw (ind) -- (pre) -- (sce);
\end{tikzpicture}
\caption{Schema scheletro della vista C.}
\label{fig:sch-C}
\end{figure}


% Bibliografia: stampa solo le voci di ref.bib citate nel testo (serve bibtex per compilare).
\newpage
\addcontentsline{toc}{section}{\refname}
\bibliographystyle{apalike}
\bibliography{ref}

\end{document}






%All other official TU Delft colours
\definecolor{donkerblauw}{RGB}{12, 35, 64}
\definecolor{turkoois}{RGB}{0, 184, 200}
\definecolor{blauw}{RGB}{0, 118, 194}
\definecolor{paars}{RGB}{111, 29, 119}
\definecolor{roze}{RGB}{239, 96, 163}
\definecolor{framboos}{RGB}{165, 0, 52}
\definecolor{rood}{RGB}{224, 60, 49}
\definecolor{oranje}{RGB}{237, 104, 66}
\definecolor{geel}{RGB}{255, 184, 28}
\definecolor{lichtgroen}{RGB}{108, 194, 74}
\definecolor{donkergroen}{RGB}{0, 155, 119}
%You can use these to change the hyperlink colour or the colour of the header or whatever. Glück Auf!
